\documentclass[conference,a4paper]{IEEEtran}
\IEEEoverridecommandlockouts
\usepackage{cite}
\usepackage{amsmath,amssymb}
\usepackage{graphicx}
\usepackage[hidelinks]{hyperref}

\begin{document}

\title{Stock Versus Flow in Supplier-Finance Disclosure:
Programme Net Flow Carries the Forecasting Content}

\author{
\IEEEauthorblockN{AVR Akshaya}
\IEEEauthorblockA{\textit{School of Business and Management}\\
\textit{CHRIST (Deemed to be University)}\\
Pune, India\\
akshaya.avr@christuniversity.in}
\and
\IEEEauthorblockN{Jerry Silvon J M}
\IEEEauthorblockA{\textit{School of Business and Management}\\
\textit{CHRIST (Deemed to be University)}\\
Pune, India\\
jerry.silvon@bbabah.christuniversity.in}
}

\maketitle

\begin{abstract}
Supplier finance lets a buyer confirm an invoice, have a bank pay the supplier
early at a discount, and settle later itself, which extends the buyer's payment
terms using the supplier's liquidity. Until recently a reader of the accounts
could not see it. ASU 2022-04 now requires a rollforward of confirmed
obligations, so the gross flows through the programme are disclosed alongside
the year-end balance. In the first mandated year 193 filers report USD 407.5 billion of obligations
added against USD 96.0 billion outstanding at the year-end, so the balance is a
thin summary of what moved. We ask whether the
flows carry information the balance does not. On 230 firm-years from 167 firms,
over a full Barth-Cram-Nelson disaggregated-accrual baseline, net programme flow
forecasts next-year operating cash flow with a coefficient of $-0.959$
($p = 0.001$), about 8.6\% of a standard deviation of the outcome. The
relationship is asymmetric: net flow survives controlling for the balance
($p = 0.025$), while the balance does not survive controlling for net flow
($p = 0.221$). Gross throughput carries nothing, whether entered alone
($p = 0.210$) or alongside the balance ($p = 0.986$), so the informative
quantity is the net direction of the programme rather than its size. We test the
result against the two mechanical explanations it invites. It is not payables
reversal: net flow is uncorrelated with the following year's change in payables
($r = 0.022$) and the coefficient is unmoved by controlling for it. It is not
year-end working-capital manipulation: the effect is present among
December-fiscal-year-end filers and stronger among the rest.
\end{abstract}

\begin{IEEEkeywords}
supply chain finance, reverse factoring, disclosure, cash flow forecasting,
working capital, ASU 2022-04, XBRL
\end{IEEEkeywords}

\section{Introduction}

In a supplier finance programme a buyer confirms that a supplier's invoice is
valid, a bank pays the supplier early at a discount, and the buyer settles with
the bank later than it would have settled with the supplier. Every party gets
something, and the buyer's working capital improvement is funded by a discount
the supplier absorbs.

For a reader of the accounts this was, for a long time, close to invisible.
Confirmed obligations sit in accounts payable rather than in debt, so a firm
could extend its terms substantially without the balance sheet showing anything
resembling borrowing. The arrangements drew regulatory attention after several
failures in which large undisclosed programmes were found to have flattered
working capital, and in September 2022 the Financial Accounting Standards Board
issued ASU 2022-04, which requires disclosure of programme obligations and,
from fiscal years beginning after 15 December 2023, a rollforward of them.

The rollforward converts a stock into a flow, which is what matters here. Across
the fiscal years beginning in 2024, the first mandated year, 193 filers report
\textbf{USD 407.5 billion} of obligations added to their programmes against
\textbf{USD 96.0 billion} outstanding at the year-end, and at the median firm one
year's additions run 4.5 times the closing balance. The balance is not measuring
scale.

That raises the question this paper answers. Does the flow carry information
about the firm's future that the balance does not? The question has a practical
edge, since the balance was disclosable a year before the rollforward was and is
the figure most data vendors carry, and it has an analytical one, because stock
and flow are different objects and the literature on supplier finance has had
access only to the stock.

The question is live rather than hypothetical. Brockbank, Hill and Mutlu
\cite{brockbank2025} examine these disclosures and find that the standard does
not improve analysts' ability to forecast cash flows. Their account of why
programmes are run also supplies the leading alternative explanation for our
result, which Section~VI-E tests.

We test it directly. On 230 firm-years from 167 filers, with next-year operating
cash flow as the outcome and the Barth, Cram and Nelson disaggregated-accrual
specification \cite{barth2001} as the baseline, net programme flow enters with a
coefficient of $-0.959$ and a cluster-robust $p$ of 0.001. A firm whose
programme expands on net, adding more than it settles, has weaker operating cash
flow the following year.

What makes it a finding rather than a correlation is the structure. Net flow
keeps its significance with the year-end balance alongside it; the balance,
marginally significant on its own, does not keep its significance with net flow
alongside it; and gross throughput is insignificant everywhere we put it. The
ordering is asymmetric and it runs against the stock.

Three contributions follow. This is the first test of the rollforward's
information content. The decomposition then shows that the informative quantity
is net direction and not gross scale, which is not what a reader of the
disclosure would guess, and we establish what the rollforward adds over
differencing two consecutive balances.

\section{Related Work}

Cash flow forecasting from accrual components is a settled empirical design.
Barth, Cram and Nelson \cite{barth2001} established that disaggregating earnings
into cash flow and the individual accruals predicts future operating cash flow
better than aggregate earnings does, building on the relation between earnings
and cash flows documented by Dechow, Kothari and Watts \cite{dechow1998}. The
accrual estimation error framework \cite{dechow2002} and the broader earnings
quality literature \cite{dechow2010} refine what the components mean, and the
predictive content of accruals shifts with the operating environment
\cite{nallareddy2020}. We adopt the specification as the baseline precisely
because it is the demanding one: a new variable that adds nothing over
disaggregated accruals adds nothing.

Supplier finance itself has been studied largely as an operations and financing
problem, covering programme mechanics and adoption \cite{pfohl2009,
gelsomino2016, wuttke2013, liebl2016}, the value of payment term extension to
the buyer \cite{lekkakos2017}, and the role of factoring in financing smaller
suppliers \cite{klapper2006}. The trade credit literature supplies the economics
of who finances whom and at what implicit cost \cite{petersen1997,
giannetti2011, klapper2012, murfin2015, chod2017}, and working capital behaviour
has its own descriptive tradition \cite{hill2010}. None of it had a disclosed,
machine-readable measure of programme activity, for the sufficient reason that
none existed.

One line of work has examined ASU 2022-04 directly. Brockbank, Hill and Mutlu
\cite{brockbank2025} identify programme users from filing text, treat supplier
financing as a binary indicator, and examine how the arrangements relate to
working-capital management, reporting incentives and analysts' cash flow
forecasts. We differ on the dimensions that matter here, in that the measure is
the continuous net flow through the programme rather than a binary indicator,
the source is the XBRL rollforward rather than filing text, and the outcome is
the firm's own realised cash flow rather than analyst forecast properties. They
are this paper's predecessor and not its rival. Their finding that programmes
are used to manage year-end payables is also the sharpest competing account of
our result, and Section~VI-E tests it.

Mandatory disclosure changes what reaches the market and at what processing cost
\cite{leuz2016, blankespoor2020, christensen2013}, and the off-balance-sheet
literature shows how arrangements that are economically financing but
presentationally something else are treated once they become visible
\cite{altamuro2014}. Supplier finance sat in exactly that position until 2022.

\section{The Disclosure}

ASC 405-50 defines a supplier finance programme by three features: the buyer
agrees with a finance provider, the buyer confirms supplier invoices as valid to
that provider, and the supplier may then request early payment from the provider
rather than from the buyer. The subtopic is disclosure-only. It does not tell a
firm whether to present the obligation as trade payable or as debt, which is
what makes a clean test of the disclosure's content possible: the standard
changed what is visible without changing the accounting underneath.

Two timing features shape any empirical design here. The balance and key terms
became effective for fiscal years beginning after 15 December 2022 and applied
retrospectively, whereas the rollforward became effective a year later, for
fiscal years beginning after 15 December 2023, and applied
\emph{prospectively}, so no retrospective activity history exists. Interim
periods require only the balance. Annual frequency is therefore not a design
choice on our part.

\section{Data and Sample}

Filings were identified by full-text search of 10-K filings for supplier finance
terminology, yielding 551 distinct filers, and structured data was then pulled
per company from the SEC XBRL company facts API. The activity elements are
\texttt{SupplierFinanceProgramObligationIncrease} and
\texttt{SupplierFinanceProgramObligationDecreaseSettlement}; the balance elements
are \texttt{SupplierFinanceProgramObligation} and its current variant.

We require both activity elements for the same annual period, since a
rollforward needs both sides, and a small number of filers tag settlements
without tagging annual additions. That yields \textbf{208 firms and 438 usable
firm-years}. Requiring a filed subsequent-year operating cash flow reduces this
to 232, and excluding financial firms following the baseline specification gives
the estimation sample of \textbf{230 firm-years from 167 firms}.

The binding constraint on sample size is filing lag rather than adoption, since
a firm-year needs both its own rollforward and the following year's cash flow,
and for fiscal years ending in late 2025 the following year has not been filed.

Table~\ref{tab:desc} reports descriptive statistics. The row that motivates the
paper is the contrast between gross and net: obligations added average 0.069 of
assets while net programme flow averages 0.0011, so the gross flows are roughly
sixty times the net at the mean. A programme moves very large volumes and nets
to very little, which is why treating gross activity and net direction as the
same variable would be a mistake.

\begin{table}[htbp]
\caption{Descriptive statistics}
\label{tab:desc}
\centering
\begin{tabular}{lrrrr}
\hline
\textbf{Variable} & \textbf{Mean} & \textbf{SD} & \textbf{Median} & \textbf{P75}\\
\hline
$\mathrm{OCF}_{t+1}$ & 0.1057 & 0.0750 & 0.0943 & 0.1352\\
$\mathrm{OCF}_{t}$ & 0.1053 & 0.0652 & 0.0934 & 0.1362\\
Programme balance & 0.0197 & 0.0392 & 0.0068 & 0.0192\\
Obligations added & 0.0694 & 0.0818 & 0.0406 & 0.0840\\
Obligations settled & 0.0682 & 0.0802 & 0.0385 & 0.0822\\
Net programme flow & 0.0011 & 0.0067 & 0.0001 & 0.0021\\
\hline
\end{tabular}
\end{table}

All variables are scaled by average total assets and winsorised at the first and
ninety-ninth percentiles.

\section{Method}

The outcome is operating cash flow in year $t+1$, scaled by average total
assets. The baseline is the Barth-Cram-Nelson specification: current operating
cash flow plus the disaggregated accruals, meaning the changes in receivables,
inventory and payables, depreciation and amortisation, and other accruals
recovered as a residual from the identity that earnings equal cash flow plus
accruals. Financial firms are excluded, as in the original. Estimation is by OLS
with standard errors clustered by firm, treated as asymptotically normal across
167 clusters.

We then add, in turn, the year-end programme balance; the two gross flows
separately; their difference as a single net term; gross throughput, defined as
the mean of additions and settlements; and combinations of these. Incremental
content is assessed by cluster-robust Wald tests of the added terms rather than
by the change in adjusted $R^2$, since the candidates are being asked to add to
a specification already explaining 76.3\% of the variance.

Including $\Delta$AP in the baseline is deliberate and consequential. Confirmed
supplier finance obligations generally sit within accounts payable, so a test
that omitted the payables accrual could recover it through the programme
variable and report a mechanical result. With $\Delta$AP controlled, net
programme flow correlates with it at only 0.372.

\section{Results}

\begin{figure*}[t]
\centering
\includegraphics[width=0.86\textwidth]{figA1.pdf}
\caption{Panel (a): obligations added during the year divided by the year-end
balance. The median firm's additions are 4.5 times its closing balance, so the
balance summarises a small fraction of the year's activity. Panel (b):
added-variable plot in standard-deviation units, with both outcome and predictor
residualised on the baseline and on the other predictor. Net programme flow
retains a negative slope against next-year operating cash flow; the year-end
balance is flat.}
\label{fig:stockflow}
\end{figure*}

\subsection{Net flow forecasts, the balance does not}

Table~\ref{tab:main} gives the nested specifications. The baseline reaches an
adjusted $R^2$ of 0.7627. Adding the year-end balance alone produces a
coefficient of $-0.124$, marginally significant at $p = 0.031$. Adding the two
gross flows separately produces coefficients that are large, opposite in sign and
close to equal in magnitude: additions enter at $-0.912$ and settlements at
$+0.883$.

\begin{table}[htbp]
\caption{Incremental content over the accrual baseline}
\label{tab:main}
\centering
\begin{tabular}{lrrr}
\hline
\textbf{Specification} & \textbf{Coef.} & \textbf{$p$} & \textbf{adj. $R^2$}\\
\hline
(1) Baseline & & & 0.7627\\
(2) $+$ Balance & $-$0.124 & 0.031 & 0.7658\\
(3) $+$ Added / Settled & $-$0.912 & 0.002 & 0.7692\\
 & $+$0.883 & 0.003 & \\
(4) $+$ Net flow & $-$0.959 & 0.001 & 0.7681\\
(5) $+$ Net flow & $-$0.776 & 0.025 & 0.7684\\
\quad\ \ and Balance & $-$0.075 & 0.221 & \\
(6) $+$ Throughput & $-$0.043 & 0.210 & 0.7637\\
(7) $+$ Throughput & $+$0.001 & 0.986 & 0.7648\\
\quad\ \ and Balance & $-$0.125 & 0.128 & \\
\hline
\end{tabular}
\end{table}

That near-symmetry is informative rather than incidental. The restriction that
the two coefficients sum to zero cannot be rejected ($\chi^2 = 0.742$,
$p = 0.389$), so the two gross flows collapse to a single net term without loss.
Estimated directly, net programme flow enters at $-0.959$ with $p = 0.001$.

The magnitude is modest, which is worth stating because a coefficient near minus
one reads as mechanical if the units are left implicit. A one standard deviation
move in net programme flow is 0.0067 of assets, so the implied effect on
next-year operating cash flow is about 0.0064 of assets: roughly 8.6\% of a
standard deviation of the outcome, or about 6\% of its mean. This is an
economically sensible signal rather than a one-for-one accounting reversal.

Specification (5) is the paper's central test, and it is deliberately the
unfavourable one, since it puts the candidate variable and the incumbent in the
same model and lets them compete. Net flow survives at $p = 0.025$. The balance
does not, at $p = 0.221$. Run the same contest with the two gross flows in place
of the net term and the balance does worse still, at $p = 0.618$.

We state what this does and does not establish. A $p$ of 0.221 is a failure to
reject, not a demonstration that the balance contains exactly nothing, and with
230 observations the test is not powerful enough to prove a null. What the
evidence establishes is an ordering: on this sample, the flow retains its
information when the stock is present and the stock does not retain its
information when the flow is present. That asymmetry is the result, and it is
the sharpest claim the data supports.

\subsection{Gross scale carries nothing}

Gross throughput is insignificant everywhere we put it, reaching only
$p = 0.210$ over the baseline and an essentially zero coefficient at
$p = 0.986$ once the balance is controlled alongside it. That is worth reporting
rather than dropping, because throughput is the proxy for programme size a
reader of the rollforward would reach for first. Its correlation of 0.72 with
the year-end balance is most of the explanation, since gross scale and the stock
carry substantially the same information and neither survives the flow.

\subsection{Forecasting outside the estimation sample}

The preceding tests ask whether net flow adds explanatory content. Anyone
intending to use the disclosure asks something narrower, which is whether it
forecasts better on observations the estimator has never encountered. We answer
that two ways, and in both the winsorising bounds derive from the training rows
alone, so the tails of held-out observations never influence estimation.

Excluding one firm at a time and predicting its firm-years from the remaining
166, root mean squared error falls from 0.03775 to 0.03732 once net programme
flow enters, and out-of-sample $R^2$ improves by 0.0058. Estimating instead
across fiscal years through 2024 and predicting 2025 reproduces that direction
on 48 held-out observations. Table~\ref{tab:oos} collects both.

\begin{table}[htbp]
\caption{Forecast accuracy outside the estimation sample}
\label{tab:oos}
\centering
\begin{tabular}{lrrr}
\hline
\textbf{Design} & \textbf{Baseline} & \textbf{$+$ Net flow} & \textbf{Change}\\
\hline
Leave-one-firm-out RMSE & 0.03775 & 0.03732 & $-$1.15\%\\
\quad out-of-sample $R^2$ & 0.7456 & 0.7514 & $+$0.0058\\
FY2025 holdout RMSE & 0.04558 & 0.04475 & $-$1.81\%\\
\quad out-of-sample $R^2$ & 0.6430 & 0.6558 & $+$0.0128\\
\hline
\end{tabular}
\end{table}

The limit of this is worth stating before anyone leans on it. A
Diebold-Mariano comparison of squared errors averaged within firm returns
$p = 0.342$, so an improvement of roughly one percent is not distinguishable
from no improvement at 167 clusters. What the exercise establishes is narrower,
and it is that the in-sample content survives the move to held-out data with its
sign intact, which a variable contributing half a point of adjusted $R^2$ could
easily have failed to do. A forecasting gain a practitioner could bank is not
established here. Demonstrating one needs the second mandated year.

\subsection{Influence, specification and competing controls}

Three separate questions about fragility turn out to share an answer, and
Table~\ref{tab:robust} reports them together.

Does a single firm carry the result? Refitting 167 times, each time excluding a
different firm, holds the coefficient between $-1.043$ and $-0.805$ and
significant at five percent in every individual fit, the weakest $p$ being
0.006. Observations prove more consequential than firms. Excluding the most
influential by DFBETA---Fluence Energy's fiscal 2024---relocates the estimate to
$-0.738$; that firm's own fiscal 2023 observation pulls oppositely, which is why
excluding Fluence altogether returns $-0.818$. Magnitude proves sensitive to a
handful of observations in a way significance does not.

Does winsorising manufacture the relationship? Without winsorising the
coefficient becomes $-1.610$, appreciably larger than the estimate we report,
and winsorising instead at the second and ninety-eighth percentiles produces
$-0.929$. The one-percent treatment is the conservative choice among the three.

Do liquidity and financing conditions explain the relationship? Introducing
cash, long-term debt and size leaves net programme flow between $-0.881$ and
$-1.110$ and significant throughout. Incomplete balance-sheet elements cost the
fully controlled specification 83 observations, and across that reduced sample
net flow already registers $-1.164$ before a single control enters, so the
restriction relocates the estimate further than the controls themselves. These
are levels rather than changes, which leaves untouched the possibility that
economic conditions move programme usage and subsequent operations together.

\begin{table}[htbp]
\caption{Net programme flow under alternative treatments}
\label{tab:robust}
\centering
\begin{tabular}{lrrr}
\hline
\textbf{Treatment} & \textbf{$N$} & \textbf{Coef.} & \textbf{$p$}\\
\hline
As reported & 230 & $-$0.959 & 0.001\\
Leave-one-firm-out, weakest of 167 & 229 & $-$0.805 & 0.006\\
Drop most influential observation & 229 & $-$0.738 & 0.013\\
Unwinsorised & 230 & $-$1.610 & 0.001\\
Fiscal-year effects & 230 & $-$1.177 & 0.001\\
$+$ cash, leverage, size & 147 & $-$1.110 & 0.032\\
\hline
\end{tabular}
\end{table}

Industry is the one cut this sample cannot support. Of the 230 firm-years, 170
are manufacturers and 41 sit in wholesale or retail trade. Within manufacturing
the estimate is $-1.748$ for SIC 30--39 against $-0.604$ for SIC 20--29, but the
interaction across that boundary returns $p = 0.185$ and interacting net flow
with payables intensity returns $p = 0.259$. We therefore report no industry
variation rather than the contrast, which is inseparable from noise at 71
observations in the smaller group.

\subsection{What the rollforward adds over differencing}

Net flow tracks the change in the disclosed balance at a correlation of 0.970,
which prompts the obvious objection that two consecutive balances would give the
same signal a year earlier. On the 209 firm-years carrying both an opening and a
closing balance the differenced balance does carry signal, at $-0.803$
($p = 0.012$) against net flow's $-1.053$ ($p = 0.002$), and the two are
collinear by construction. The rollforward's advantage is availability rather
than strength: it yields the flow from a single filing, and 9.1\% of our
firm-years have no prior-year balance at all, which makes differencing there not
inconvenient but impossible.

\subsection{Two mechanical explanations, tested}

A coefficient close to minus one on a working-capital variable invites the
suspicion that an accounting identity is unwinding rather than a forecast being
made. Two versions of that suspicion are worth taking seriously, and both are
testable.

The first is payables reversal. Net flow in year $t$ moves accounts payable at
the year-$t$ balance sheet date, that balance is the opening position for the
change in payables during $t+1$, and that change sits inside the outcome. The
baseline controls $\Delta$AP in year $t$, which handles the contemporaneous
channel, but not the one that runs through the following year. We recovered
year-$t{+}1$ payables from company facts for 200 of the 230 firm-years and tested
it directly. Net programme flow correlates with the following year's change in
payables at $r = 0.022$. On that subsample the coefficient is $-0.885$
($p = 0.006$); adding $\Delta$AP$_{t+1}$ as a control leaves it at $-0.892$
($p = 0.007$); and removing the payables change from the outcome entirely leaves
it at $-0.950$, with the standard error rising from 0.325 to 0.508 and
significance falling to $p = 0.062$. The point estimate does not move when the
payables channel is closed. What the identity story predicts is not what the
data shows.

The second is year-end manipulation. If programmes are used to inflate payables
at the reporting date, as Brockbank et al. \cite{brockbank2025} document, then
cash flow in $t$ is flattered and cash flow in $t+1$ is depressed, and our
coefficient follows with no supply-chain economics involved. Their evidence
locates that behaviour among firms with December fiscal year-ends, which gives a
direct test. Splitting on year-end month, net flow enters at $-0.641$
($p = 0.038$) among the 156 December firm-years and at $-1.868$ ($p = 0.015$)
among the 74 others, and the interaction between net flow and a December
indicator is insignificant ($p = 0.252$). The effect is present in both groups
and larger outside the group where the manipulation account predicts it should
concentrate.

Neither test proves a mechanism. Together they rule out the two explanations a
reader reaches for first.

\subsection{Early adopters}

Since the rollforward became compulsory only at the boundary described in
Section~III, part of our sample is voluntary early disclosure, and firms that
volunteer plausibly differ from those that wait. Splitting on that boundary,
net flow enters at $-0.668$ on the 157 mandated firm-years and at $-1.195$ on
the 73 voluntary ones.

The mandated subsample alone does not reject zero ($p = 0.163$), with a standard
error of 0.480 against 0.299 pooled. We report that plainly. The two estimates
are not statistically distinguishable from one another ($z = 0.79$,
$p = 0.431$), so the split gives no evidence that the pooled result is an
artifact of early adopters; what it shows is that 157 single-year observations
are too few to establish the effect on their own. A reader who wants it
demonstrated on mandated filings alone will have to wait for the second mandated
year.

\section{Discussion}

The economics behind the sign are not mysterious. A programme expanding on net,
adding more confirmed obligations than it settles, is extracting more liquidity
from its supply chain than it returns. That is a financing inflow presented
inside working capital, and the finding is that it reverses, roughly a year
before the cash flow arrives.

The stock-versus-flow distinction matters to anyone building on this data. The
year-end balance is a point observation of a quantity that turns over several
times a year, and for the median firm in our sample it captures less than a
quarter of what passed through, so two firms reporting identical balances may
have run programmes of very different size and of opposite direction. Vendors
and researchers carrying the balance because it arrived first are carrying the
less informative of the two disclosed figures.

For a user of the accounts the practical instruction is specific. An analyst
building a cash flow forecast should read additions and settlements rather than
the closing obligation, and should read their difference rather than either
alone, because a firm whose programme is expanding is reporting operating cash
flow that its own supply chain is funding. The magnitude is modest. A standard
deviation of net flow maps to roughly 8.6 percent of a standard deviation of
next year's cash flow, which tilts a forecast rather than replacing it. A credit
analyst has the sharper use, since net flow gives the direction of a financing
line that is unpriced, undated and absent from debt---the three properties that
make it worth watching. For the disclosing firm the
implication runs the other way, in that the flows now reveal what the balance
concealed, and a programme expanded to improve a year-end figure is visible in
the rollforward of the year it is expanded.

There is a standard-setting reading as well. ASU 2022-04's stated objective is to
let users understand the nature, activity and potential magnitude of these
programmes, and the rollforward was the element that made activity visible. On
this evidence the rollforward is the part of the disclosure that does the work,
and the balance, which arrived a year earlier and imposes a lighter burden, is
the part that does not add to it. That ordering is worth knowing before anyone
proposes simplifying the requirement.

\subsection{Limitations}

The sample is one country, one standard and the disclosure's first years, and
230 firm-years is small for this literature. Firm fixed effects are not
estimable at this panel depth. The result is a forecasting relationship and we
do not claim a causal channel from programme expansion to subsequent cash flow,
since firms choose how much programme to run and the choice may respond to
conditions that also drive next year's operations. Confirmed obligations are not
necessarily financed obligations, because a supplier need not take the early
payment, and the US disclosure does not separate the two.

Generalisation beyond US filers is not a matter of effort. IAS 7 paragraph 44H
requires the carrying amounts at the beginning and end of the period and the
type and effect of non-cash changes in them, with no equivalent of the added and
settled lines, so the quantity measured in this paper does not exist in IFRS
disclosure and a cross-country test of it cannot presently be constructed. The
obvious extension can be. A second mandated year roughly doubles the panel and
brings firm fixed effects and a mandate-only test within reach, which is what
the result now most needs, and whether market participants price the flow
becomes answerable once forecast and rating data are matched to the rollforward.

\section{Conclusion}

ASU 2022-04 made the flows through supplier finance programmes visible for the
first time, and in the first mandated year 193 filers put more than four hundred
billion dollars of confirmed obligations through arrangements whose year-end
footprint was a quarter of that. The flows do say something the balance does
not. Over a full disaggregated-accrual baseline net programme flow forecasts
next-year operating cash flow at $-0.959$, net flow survives the presence of the
balance while the balance does not survive the presence of net flow, and gross
throughput carries nothing.

The quantity that matters is the direction a programme is moving, not its size,
and until the rollforward was required it could not be read from the accounts at
all.

\section*{Data and Reproducibility}

All filings are public. Panel construction, estimation and exhibit scripts, with
the result artifacts behind every figure reported here, are at
\url{https://github.com/jerrysilvon/icisas2026-replication}; the analysis makes
no network calls once company facts are cached.

\begin{thebibliography}{00}

\bibitem{barth2001} M. E. Barth, D. P. Cram, and K. K. Nelson, ``Accruals and the prediction of future cash flows,'' \textit{Account. Rev.}, vol. 76, no. 1, pp. 27--58, 2001.

\bibitem{dechow1998} P. M. Dechow, S. P. Kothari, and R. L. Watts, ``The relation between earnings and cash flows,'' \textit{J. Account. Econ.}, vol. 25, no. 2, pp. 133--168, 1998.

\bibitem{dechow2002} P. M. Dechow and I. D. Dichev, ``The quality of accruals and earnings: The role of accrual estimation errors,'' \textit{Account. Rev.}, vol. 77, no. s-1, pp. 35--59, 2002.

\bibitem{dechow2010} P. Dechow, W. Ge, and C. Schrand, ``Understanding earnings quality: A review of the proxies, their determinants and their consequences,'' \textit{J. Account. Econ.}, vol. 50, no. 2--3, pp. 344--401, 2010.

\bibitem{nallareddy2020} S. Nallareddy, M. Sethuraman, and M. Venkatachalam, ``Changes in accrual properties and operating environment: Implications for cash flow predictability,'' \textit{J. Account. Econ.}, vol. 69, no. 2--3, p. 101313, 2020.

\bibitem{pfohl2009} H. Pfohl and M. Gomm, ``Supply chain finance: Optimizing financial flows in supply chains,'' \textit{Logist. Res.}, vol. 1, no. 3--4, pp. 149--161, 2009.

\bibitem{gelsomino2016} L. M. Gelsomino, R. Mangiaracina, A. Perego, and A. Tumino, ``Supply chain finance: A literature review,'' \textit{Int. J. Phys. Distrib. Logist. Manag.}, vol. 46, no. 4, 2016.

\bibitem{wuttke2013} D. A. Wuttke, C. Blome, K. Foerstl, and M. Henke, ``Managing the innovation adoption of supply chain finance: Empirical evidence from six European case studies,'' \textit{J. Bus. Logist.}, vol. 34, no. 2, pp. 148--166, 2013.

\bibitem{liebl2016} J. Liebl, E. Hartmann, and E. Feisel, ``Reverse factoring in the supply chain: Objectives, antecedents and implementation barriers,'' \textit{Int. J. Phys. Distrib. Logist. Manag.}, vol. 46, no. 4, 2016.

\bibitem{lekkakos2017} S. D. Lekkakos and A. Serrano, ``Reverse factoring: A theory on the value of payment terms extension,'' \textit{Found. Trends Technol. Inf. Oper. Manag.}, vol. 10, no. 3--4, pp. 270--288, 2017.

\bibitem{klapper2006} L. Klapper, ``The role of factoring for financing small and medium enterprises,'' \textit{J. Bank. Finance}, vol. 30, no. 11, pp. 3111--3130, 2006.

\bibitem{petersen1997} M. A. Petersen and R. G. Rajan, ``Trade credit: Theories and evidence,'' \textit{Rev. Financ. Stud.}, vol. 10, no. 3, pp. 661--691, 1997.

\bibitem{giannetti2011} M. Giannetti, M. Burkart, and T. Ellingsen, ``What you sell is what you lend? Explaining trade credit contracts,'' \textit{Rev. Financ. Stud.}, vol. 24, no. 4, pp. 1261--1298, 2011.

\bibitem{klapper2012} L. Klapper, L. Laeven, and R. Rajan, ``Trade credit contracts,'' \textit{Rev. Financ. Stud.}, vol. 25, no. 3, pp. 838--867, 2012.

\bibitem{murfin2015} J. Murfin and K. Njoroge, ``The implicit costs of trade credit borrowing by large firms,'' \textit{Rev. Financ. Stud.}, vol. 28, no. 1, pp. 112--145, 2015.

\bibitem{chod2017} J. Chod, ``Inventory, risk shifting, and trade credit,'' \textit{Manag. Sci.}, vol. 63, no. 10, pp. 3207--3225, 2017.

\bibitem{hill2010} M. D. Hill, G. W. Kelly, and M. J. Highfield, ``Net operating working capital behavior: A first look,'' \textit{Financ. Manag.}, vol. 39, no. 2, pp. 783--805, 2010.

\bibitem{leuz2016} C. Leuz and P. D. Wysocki, ``The economics of disclosure and financial reporting regulation: Evidence and suggestions for future research,'' \textit{J. Account. Res.}, vol. 54, no. 2, pp. 525--622, 2016.

\bibitem{blankespoor2020} E. Blankespoor, E. deHaan, and I. Marinovic, ``Disclosure processing costs, investors' information choice, and equity market outcomes: A review,'' \textit{J. Account. Econ.}, vol. 70, no. 2--3, p. 101344, 2020.

\bibitem{christensen2013} H. B. Christensen, L. Hail, and C. Leuz, ``Mandatory IFRS reporting and changes in enforcement,'' \textit{J. Account. Econ.}, vol. 56, no. 2--3, pp. 147--177, 2013.

\bibitem{altamuro2014} J. Altamuro, R. Johnston, S. Pandit, and H. Zhang, ``Operating leases and credit assessments,'' \textit{Contemp. Account. Res.}, vol. 31, no. 2, pp. 551--580, 2014.

\bibitem{brockbank2025} B. Brockbank, M. Hill, and S. Mutlu, ``Supplier financing: Incentives, manipulation, and disclosure effects,'' working paper, Dec. 2025. [Online]. Available: https://ssrn.com/abstract=5933157

\end{thebibliography}

\end{document}
